\chapter{Background}
\label{sec:background}

This chapter briefly gives introductory descriptions of blockchain, smart
contracts, decentralized finance, and transaction execution.


\section{Blockchain}


Blockchain is a distributed ledger that broadcasts and stores information of
transactions across different parties. 
A blockchain consists of a growing
number of blocks and a consensus algorithm which decides the order of blocks.
Each block is constituted of transactions.
%While Satoshi Nakamoto invented the first decentralized cryptocurrency Bitcoin
%in 2008~\cite{nakamoto2008bitcoin}, the 
Ethereum~\cite{ethereum22,buterin2014ethereum} is the first blockchain to
support, store, and execute Turing complete programs, known as smart contracts.
In particular, the Ethereum blockchain is constituted of a global state and
transactions that modify the state. 
%Ethereum supports two types of accounts:
%user accounts and smart contract accounts, and each account is associated with
%a unique address.  
New popular blockchains
%, e.g., Binance Smart
%Chain~\cite{bsc},  Avalanche~\cite{avalanche},
%Conflux~\cite{li2020decentralized}, and Near~\cite{near}, Fantom~\cite{fantom}
work with Ethereum virtual machine (EVM) based execution environment due to the
popularity of the EVM stack for developers. 

%Smart contracts can create internal transactions to other smart contracts. The internal transactions are nested from a top-level transaction, an external transaction, initiated by a user account. If a transaction is aborted then the effects of all nested internal transactions will be reversed. 




\section{Smart Contracts} 
%Smart contracts are programs created to automatically execute those
%transactions. 
Each smart contract is associated with a unique address, a
persistent account's storage trie, a balance of native tokens, e.g., Ether in
Ethereum and BNB in Binance Smart Chain, and bytecode (e.g., EVM
bytecode~\cite{ethereum22,buterin2014ethereum}) that executes incoming
transactions to change the storage and balance. Users interact with a smart
contract by issuing transactions from their user accounts to the contract
address. 
%Smart contracts become
%immutable once deployed to the blockchain.  
Smart contracts can also interact with other smart contracts as function calls.
%create internal
%transactions to interact with other smart contracts. The internal transactions
%are nested from a top-level transaction, an external transaction, initiated by
%a user account. If an external transaction is aborted then the effects of all
%nested internal transactions will be reversed as well. 
Currently, there are several human-readable high-level programming languages,
e.g., Solidity~\cite{solidity} and Vyper~\cite{vyper}, to write smart contracts
that compile to the EVM bytecode. 
%In Figure~\ref{fig:ERC20}. we show an
%excerpt of the ERC20 token smart contract written in Solidity. The listed
%function {transfer} allows an address, {from}, to transfer ERC20 tokens to
%another address, {to}. 

\input{solidity-highlighting.tex}
\begin{figure}[t]
    \centering
    \begin{minipage}[h]{0.9\linewidth}
    \lstinputlisting{codes/ERC20.sol}
    \end{minipage}
    \caption{An excerpt of ERC20 smart contract.}
    \label{fig:ERC20}
    \vspace{-2mm}
\end{figure}




\section{Decentralized Finance (DeFi)}

Decentralized finance (DeFi) is a peer-to-peer financial ecosystem built on top of
blockchains~\cite{wust2018you}. The building blocks of DeFi are smart contracts
that manages digital assets. 
%that constitute the DeFi
%protocols~\cite{popescu2020decentralized,cao2021flashot}. 
A few DeFi protocols dominate the DeFi market and serve as references for other
decentralized applications: stable coins (e.g., USDC and USDT), price oracles,
decentralized exchanges, and lending and borrowing platforms. 
%\subsubsection{Flash loan} 
In DeFi, a special type of loan called \emph{flash
loan} allows lenders to offer loans to borrowers without upfront collaterals
deposits. The loan is only valid within a single transaction and must be repaid
with fees before the completion of the transaction (otherwise the transaction
is reverted). 
 
%These money legos can be composed to form various applications. For instance, 1inch~\cite{1inch} is a DEX aggregator that assist users to find the best crypto prices across DEXes, and reroutes trades between them. 
%Composability is one of the most important feature of DeFi.
%However, referencing external protocols also introduces new risks to a DeFi application, since external protocols' unusual behaviors could easily lead to the applications' unexpected behaviors without rigorous checks. 



\section{Transaction Execution} 

A transaction is constituted of a sender address, a recipient address, a
transferred value of native token (can be zero), a transaction data (can be
empty), and a gas value to pay the transaction fees. If the transaction
recipient address is associated with a user account then the transaction is
called \emph{simple payment} transaction that transfers the value of native
token from the sender account to the recipient account and the transaction's
data is empty in this case. Otherwise, if the transaction recipient address is
associated with a smart contract account then the transaction's data identifies
a method of the recipient smart contract's code together with arguments passed
to execute the method. When the transaction is received by the block miners,
the corresponding smart contract's method is executed in the EVM, modifying the
smart contract's balance and storage trie accordingly. Note that the execution
of each EVM instruction, such as read/write operations on the underlying
storage trie, is associated with a gas fee. The transaction's gas value must
exceed the accumulated gas fees at the end of execution, otherwise the
transaction is aborted~\cite{ethereum22,buterin2014ethereum}.



